\section{The High Priestess}


When God manifested in the created plane, the Two inescapably arose. First of all, as secondary to God, and secondly as articulations within itself, to create the pyramid and hierarchy of all Being. There is nothing wrong with Two: as Cologero notes, the animus within the woman is necessary for her healthy development. The Two is very good in itself, but not by itself. This is noted earlier on as \emph{it is not good for man to be alone}. This would also apply to the feminine part within man's own interior.

Watch a political debate, read a cutting edge expose of the Gospels, pick up a history book, or (worse) sample a modern textbook. For that matter, pick an acquaintance's brain. What you will inevitably almost always discover is someone for whom the Two is not merely a good, but is actually God, subsisting alone. Not only have they made a false god out of what is created, but (for them) the One and indeed any other competing Twos simply don't exist. Thus, those who think the Catholic Church was part of a vast cosmic conspiracy involving the Demi-Urge suppressing channellers from the stars simply never get around to asking themselves the simple question, If God is a fake God (\textbf{Yahweh is the Demi-urge}) then how do I know the One True Light God is not just in on the fix? What if the deception \emph{is even more elaborate} than I first believed? No for them what is an important psychological insight gets distorted into a cosmology, an ideology, and at last, a madness. They cannot square the circle. This phenomena is common today.

Psychologically and humanly speaking, people like modern liberals, civil rights activists, free thinkers and others may certainly possess a modicum of truth. Yet, what lie does not contain modicums of truth? It is not hard to be partly right. But this is not what is claimed: it is claimed that this is the truth, the whole truth, nothing but the truth, and (most importantly) the only truth which could possibly matter. And it must be taken \emph{as presented}, which is to say, in isolation from the One. It is, at root, a narrative created by the demonic power of the egocentric mind, which attempts to enslave others to its dominating will. This is the opposite of the what the Priestess was created to be: instead of the \emph{alma mater}, we have the devouring wrath of the psychotic Lilith, who hates Eve, and attempts to destroy Adam (this is how salvation comes from the woman, for Eve will oppose Lilith).

We see at work here the very power of distortion and twisting which they claim to see in the One. To quote Phillip Dick, ``\textbf{Those who fight the Empire, become the Empire}''. Or to quote Jesus, ``resist not evil, but overcome evil with Good''. Anyone who uses the Ring of Sauron against Sauron is doomed to either become like him, or to be turned into a slave. God Himself does not overcome evil with evil, but rather overturns Hell from the inside out. Good triumphs because it remains unchanged when Evil touches it, and is capable of changing Evil into its opposite. Evil cannot make this claim, and therefore, has no true power. Evil can only uncover what is already Evil, and test what is Good.

This is not to say one cannot be a warrior and a man of the One. There are beautiful stories about how to reconcile the two, \& many men have done it extremely well, so well they merit great titles and honors, and perhaps even more than man can say. Giorgio\footnote{\url{https://www.gornahoor.net/?p=6036}} has pointed out that the warrior (however) is female in relation to the brahmin or priest. That is, the warrior has to overcome the flesh to arrive at the one, he has to fight himself. He has to struggle against the tendency to carve out a separate kingdom away from the Law of the One. This is the first temptation faced by man, to divide himself into Two and to make a kingdom of his own, separate from the aboriginality of the Self Beyond the Selves.

How can this be? How can the virile, manly warrior be female? When the universe is taken hierarchically, we can say that that which is above something is male in relation to it, and that which is below something is female in relation to it. Things below revolve around, cling to, and require for their health and sanity the relation to that which stems from above. When something is detached, it careens off into the outer darkness,\footnote{\url{https://biblehub.com/jude/1-13.htm}} and although it may retain its own hierarchy for a time, there is a process of decay that sets in immediately. The horizontal cannot sustain itself by itself. It requires a link to the vertical. Thus, everything is male/female, and yet we can properly speak of certain things as ``male'' or ``female'' in configuration to each other.

Neither is privileged. Without the female, the male does not make itself known. If there is a privilege, it consists of duties. To quote the Baron, \textbf{loving,proud obedience and humble command} are the duties of respective stations. Since all of Being involves climbing the ladder or the Tree, there is no Black Iron Prison\footnote{\url{https://en.wikipedia.org/wiki/VALIS}} of compartments or boxes, but rather a kind of dance. If you are living in a Black Iron Prison, it is because your soul (to that extent) is tainted and evil, and God manifests to you as the oppressor. The eye, being not single, cannot see the Light. The psychological insight (however) of this state is correct: you are seeing that the God you made in your own image is actually a Demi-Urge. Angels appear as demons, men of good breeding and sound sense are tyrants, Western civilization is the greatest evil on earth, etc. The only way out is ``deeper in'', that is to say, to pull the beam out of your own Eye.

The lesson of the second Card is that \textbf{Two cannot exist alone}. Which (put this way) is irrefutable. QED. Or, put another way, \textbf{Two can only exist together as One}.

Tomberg gives an example of what happens when the Two wants to be alone:

\begin{quotex}
Passing on to mysticism which has not given birth to gnosis, magic and Hermetic philosophy -- such a mysticism must, sooner or later, necessarily degenerate into `spiritual enjoyment' or `intoxication'. The mystic who wants only the experience of mystical states without understanding them, without drawing practical conclusions from them for life, and without wanting to be useful to others, who forgets everyone and everything in order to enjoy the mystical experience, can only be compared to a spiritual drunkard.
\end{quotex}


He mentions that those who develop the spiritual sense develop a ``touch'' which allows him to apply and understand what he has experienced. This is analogous to the virginal pool, which more deeply reflects (until it engenders or creates) what is shining on it from above : the Sun gives birth over the void, or the waters. One begins to examine the thoughts that enters one's mind, to make them pass ``tests'' (smell, touch, hearing) before true gnosis (or sight) is finally engendered, much as a baby learns to see in three dimensions and to arrange shapes in proper contours to the forms emanating them. This is the meaning of ``become as little children'', who are female/passive at first as regards that which is already fully grown and solar.

The Gnostic challenge to the Church is an opportunity to recover what was Once: an exoteric Church which is sustained by a Life, to reunify the split between Gnosis and Faith that spiraled out of control in the early development of the West. Neither the New Age or Gnostics on the one hand, nor the purely exoteric believer on the other can long subsist in division. Such a division engenders monsters, more confusion, and less and less life. We can apply this pattern to the post-colonial schools of thought in literature, to the revolutionary and reactionary schools of thoughts in politics, and to the endless culture debates in sociology and geopolitics.

\begin{quotex}
``On the part of the human being it is the act of daring to aspire to the supreme Reality, and this act is real and effective only when the soul is serene and the body completely relaxed -- without smoke and crackling fire.''
\end{quotex}


We remember the One, because gnosis and magic and alchemy ``do not dazzle God.'' In the eyes of God, their practitioners are ``dear sheep to him: in his consideration of them he desires that they shall never go astray and that they shall have life increasingly and unceasingly''.

The yoke is easy, the burden is light. There are more dangers ahead, but to learn the lesson of the One is to pass into abundant Life.


\flrightit{Posted on 2014-05-10 by Logres}

